\section{Docker Compose}

\begin{frame}
  \frametitle{O Problema com Múltiplos Containers}
  Aplicações modernas envolvem vários serviços: Banco de dados, cache, filas, etc.

  \vspace{0.5em}
  \begin{block}{Boa Prática}
    Cada container deve fazer uma única função e fazê-la bem.
  \end{block}

  \vspace{0.5em}
  Usar vários \texttt{docker run} exige gerenciar manualmente networks
  e diversas flags de configuração.
\end{frame}

\begin{frame}
  \frametitle{O que é o Docker Compose?}
  Define todos os containers e suas configurações em um único arquivo YAML.

  \vspace{0.5em}
  \begin{itemize}
    \item Qualquer pessoa que clone o repositório sobe toda a aplicação com um único comando
    \item Networks, volumes e configurações declaradas no próprio arquivo
  \end{itemize}

  \vspace{0.5em}
  \begin{columns}
    \column{0.5\textwidth}
      \textbf{Dockerfile}
      \begin{itemize}
        \item Instrui como \textbf{construir} uma imagem
      \end{itemize}

    \column{0.5\textwidth}
      \textbf{Compose file}
      \begin{itemize}
        \item Define quais containers \textbf{executar} e como configurá-los
      \end{itemize}
  \end{columns}
\end{frame}

\begin{frame}[fragile]
  \frametitle{Estrutura do compose.yaml}
  \small
\begin{verbatim}
services:
  nome-do-servico:
    image: imagem:tag
    build: ./caminho
    ports:
      - "8080:80"
    environment:
      - VARIAVEL=valor
    volumes:
      - meus-dados:/app/data
    networks:
      - minha-rede
    depends_on:
      - outro-servico

volumes:
  meus-dados:

networks:
  minha-rede:
\end{verbatim}
\end{frame}

\begin{frame}
  \frametitle{Campos Principais do compose.yaml}
  \begin{itemize}
    \item \texttt{services}: define cada container da aplicação
    \item \texttt{image}: imagem Docker a ser usada
    \item \texttt{build}: builda a imagem a partir de um Dockerfile local
    \item \texttt{ports}: mapeamento \texttt{HOST:CONTAINER}
    \item \texttt{environment}: variáveis de ambiente
    \item \texttt{volumes}: montagem de volumes ou bind mounts
    \item \texttt{networks}: redes às quais o serviço se conecta
    \item \texttt{depends\_on}: ordem de inicialização entre serviços
  \end{itemize}
\end{frame}

\begin{frame}
  \frametitle{Comandos Básicos}
  Subir todos os serviços:

  \vspace{0.3em}
  \texttt{docker compose up -d --build}

  \vspace{0.5em}
  Derrubar todos os serviços:

  \vspace{0.3em}
  \texttt{docker compose down}

  \vspace{0.5em}
  Ver status e logs:

  \vspace{0.3em}
  \texttt{docker compose ps}

  \vspace{0.3em}
  \texttt{docker compose logs -f <servico>}
\end{frame}
